\documentclass{prosper}
\usepackage{amsmath,amsfonts,amsthm}
\usepackage{epsfig}
\include{Qcircuit}

\renewcommand{\>}{\rangle}
\newcommand{\<}{\langle}

\newcommand{\sig}{\sigma}

\newcommand{\C}{{\mathbb C}}

\def\pmat#1{\begin{pmatrix} #1 \end{pmatrix}}
%%%%%%%%%%%%%%%%%%%%%%%%%%%%%%%%%%%%%%%%%%%%%%%%%%%%%%%%%%%%%%%%%%%%%%%%%%
% Macros for drawing braids
%%%%%%%%%%%%%%%%%%%%%%%%%%%%%%%%%%%%%%%%%%%%%%%%%%%%%%%%%%%%%%%%%%%%%%%%%%

\def\r{\vtwist, c+(2,1)}
\def\l{\vtwistneg, c+(2,1)}
\def\s{\POS c+(0,0) *{};p-(0,1) *{} **\dir{-};c+(1,0)}
\def\N{\POS "A"-(0,1)="A";"A"}
\def\braid#1#2{\,\,\xy 0;<.4cm,0cm>:<0cm,.5cm>:: \POS (0,#1)="A";"A" #2 \endxy\,\,}
\def\caps#1#2{\POS c="B";"B"+(#1,0),\vcap[#2];"B";"B"}
\def\cups#1#2{\POS c="B";"B"+(#1,0),\vcap[-#2];"B";"B"}

%%%%%%%%%%%%%%%%%%%%%%%%%%%%%%%%%%%%%%%%%%%%%%%%%%%%%%%%%%%%

\usepackage{youngtab}

\title{The Jones Polynomial:\\ Quantum Algorithms \& Applications in\\ Quantum Complexity Theory}
%\subtitle
\author{Pawel M. Wocjan\\[1ex]School of Electrical Engineering and Computer Science\\ University of Central Florida}
\email{wocjan@cs.ucf.edu}

\begin{document}

\maketitle

\begin{slide}{Organization of the talk}
\begin{itemize}
\item Short introduction to quantum computing
\item Links as closures of braids
\item Jones polynomial
\item Jones-Wenzl representations of the braid group
\item Representation-theoretic formula for the Jones polynomial
\item Quantum computation as approximate evaluation of the Jones polynomial
\end{itemize}
\end{slide}


\begin{slide}{Quantum computational power}
in words:
\begin{itemize}
\item the quantum computer is very efficient at multiplying certain big sparse unitary matrices and estimating entries of the resulting product
\item this is the hardest problem which can be efficiently solved by a quantum computer
\end{itemize}

\bigskip
everything is made more precise by the quantum circuit model
\end{slide}

\begin{slide}{Quantum circuit model}
\begin{itemize}
\item let $\mathcal{H}$ be the Hilbert space given by the $n$-fold tensor product of $\mathbb{C}^2$ 
\[
\mathbb{C}^2\otimes\mathbb{C}^2\otimes\cdots\otimes\mathbb{C}^2
\]
and $SU(2^n)$ be the special unitary group acting on $\mathcal{H}$
\item $SU(2^n)$ is generated by embeddings of unitary matrices acting on $\mathbb{C}^2$ (single qubits) and $\mathbb{C}^2\otimes\mathbb{C}^2$ (two qubits); these are refered to as single qubit and two qubit \emph{gates}
\item the \emph{complexity} $\kappa(U)$ of $U\in SU(2^n)$ is the minimal number of gates such that their product gives $U$
\item the complexity $\kappa_\epsilon(U)$ is the minimal number of gates such their product yields $V$ and 
\[
\| U - V \| \le \epsilon
\]
\end{itemize}

\end{slide}

%%%%%%%%%%%%%%%%%%%%%%%%%%%%%%%%%%%%%%%%%%%%%%%%%%%%%%%%%%%%%%%%%%%%%%%%%%%%%%%%%%%%%%%%%%%

\begin{slide}{Quantum circuit model}
\begin{center}
\Qcircuit @C=0.7em @R=0.7em {
\lstick{} & \gate{U_{1\,}}     & \multigate{1}{U_4} & \qw                & \qw \\
\lstick{} & \multigate{1}{U_2} & \ghost{U_4}        & \gate{U_6}         & \qw \\
\lstick{} & \ghost{U_2}        & \qw                & \multigate{1}{U_7} & \qw \\
\lstick{} & \multigate{1}{U_3} & \gate{U_5}         & \ghost{U_7}        & \qw \\
\lstick{} & \ghost{U_3}        & \qw                & \qw                & \qw
}
\end{center}

\[
(I_2 \otimes U_6 \otimes U_7 \otimes I_2)\cdot
(U_4\otimes I_2 \otimes U_5 \otimes I_2)\cdot(U_1 \otimes U_2 \otimes U_3)
\]
\end{slide}

%%%%%%%%%%%%%%%%%%%%%%%%%%%%%%%%%%%%%%%%%%%%%%%%%%%%%%%%%%%%%%%%%%%%%%%%%%%%%%%%%%%%%%%%%%%%

\begin{slide}{Universal set \& Solovay-Kitaev theorem}

a \textcolor{green}{universal set ${\cal G}$} is a finite
subset of $SU(4)$ (closed under taking inverses) such that the subgroup $\langle {\cal G} \rangle$
generated by ${\cal G}$ is dense in $SU(4)$

\bigskip
\textcolor{green}{Solovay-Kitaev Theorem}:\\
let ${\cal G}$ be a universal set and $\delta>0$

\medskip
an arbitrary $U\in SU(4)$ can be approximated to within a distance
$\delta$ by a sequence of ${\rm poly log}(1/\delta)$ gates
from ${\cal G}$

\medskip
moreover, such a sequence be ``compiled'' by a classical computer in
${\rm poly log}(1/\delta)$ time

\medskip
$\Longrightarrow$ \textcolor{green}{the definition of complexity does not depend ``too much'' on the particular choice of ${\cal G}$}
\end{slide}

%%%%%%%%%%%%%%%%%%%%%%%%%%%%%%%%%%%%%%%%%%%%%%%%%%%%%%%%%%%%%%%%%%%%%%%%%%%%%%%%%%%%%%%%%%%%%%

\begin{slide}{Approximate Circuit Matrix Entry}
\begin{itemize}
\item let $U$ be a unitary matrix given by a quantum circuit on $n$ qubits with $m=poly(n)$ gates
and $\epsilon=1/poly(n)$
\item estimate the first diagonal entry of $U$ with precision $\epsilon$, that is, output $u$ such that
\[
|u-\<00\ldots 0|U|00\cdots 0\>|\le \epsilon
\]
\end{itemize}
there is a quantum algorithm with running time which is only polynomial in $n$

\medskip
the best classical algorithms have exponential running time

\medskip \textcolor{red}{ACME is a complete problem for the quantum complexity class Bounded error Quantum Polynomial time (BQP)}
\end{slide}

%%%%%%%%%%%%%%%%%%%%%%%%%%%%%%%%%%%%%%%%%%%%%%%%%%%%%%%%%%%%%%%%%%%%%%%%%%%%%%%%%%%%%%%%%%%%

\begin{slide}{Quantum circuit model}
the quantum circuit model is a mathematical specification of what (we hope) quantum computer could do 

\medskip
it is justified by the fundamental laws of quantum mechanics

\medskip
there are many different proposals for quantum hardware which implement this model
\begin{itemize}
\item ion traps
\item optical lattices
\item quantum dots
\item nuclear magnetic resonance
\item ...
\end{itemize}
a different model of computation is given by so-called \textcolor{blue}{topological quantum computation}
\end{slide}

%%%%%%%%%%%%%%%%%%%%%%%%%%%%%%%%%%%%%%%%%%%%%%%%%%%%%%%%%%%%%%%%%%%%%%%%%

\begin{slide}{Topological Quantum Computing}
particle exchanges in a 2+1 dimensional quantum system of $n$
identical particles are governed by the braid group $B_n$

\medskip
\quad\epsfxsize0.75\textwidth\epsfbox{braiding.eps}

particles move along the arrows, so that the original particle
positions are once again occupied

\medskip
the braid describes their trajectories in space-time

\end{slide}

%%%%%%%%%%%%%%%%%%%%%%%%%%%%%%%%%%%%%%%%%%%%%%%%%%%%%%%%%%%%%%%%%%%%%%%%%

\begin{slide}{What happens when we braid particles?}

depends on the type of particles
\begin{itemize}
\item bosons \quad\,\, $\rightarrow$\quad nothing 

\medskip
\item fermions \quad$\rightarrow$\quad wavefunction acquires $\pm 1$ factor

\medskip 
\item abelian anyons\hspace{1.3cm} $\rightarrow$\quad phase factor $e^{2\pi i\varphi}$

\medskip
\item \textcolor{green}{non-abelian anyons} \quad$\rightarrow$\quad
\textcolor{red}{nontrivial transformation}

\medskip
(quantum Hall quasi-particles and quasi-holes)
\end{itemize}

\medskip
\textcolor{blue}{certain anyons could be used for realizing universal
quantum computation with built-in fault-tolerance} 

\bigskip
fault-tolerance due to topological properties

\end{slide}

%%%%%%%%%%%%%%%%%%%%%%%%%%%%%%%%%%%%%%%%%%%%%%%%%%%%%%%%%%%%%%%%%%%%%%%%%

\begin{slide}{Mathematical description}

\textcolor{green}{state space}: some Hilbert space $V$

\bigskip
\textcolor{green}{elementary operations}: generators $\sigma_i^{\pm
1}\in B_n$ 

\bigskip

\bigskip
let $\ket{\psi}\in V$ be the state of the topological quantum computer

\bigskip 
how does \textcolor{green}{$\ket{\psi}$ change} if we apply
the elementary operation $\sigma_i^{\pm 1}$?
\[
\ket{\psi}\mapsto {\red \pi(\sigma_i^{\pm 1})}\ket{\psi}
\]

\medskip
where \textcolor{red}{$\pi$ is a unitary representation} of the braid
group $B_n$
\end{slide}



%%%%%%%%%%%%%%%%%%%%%%%%%%%%%%%%%%%%%%%%%%%%%%%%%%%%%%%%%%%%%%%%%%%%%%%%%%%%%%%%%%%%%%%%%%%%



\begin{slide}{Braid group $B_n$} 

the braid group $B_n$ on $n$ strands is generated by
$\sig_1,\sig_2,\dotsc,\sig_{n-1}$ satisfying
\begin{eqnarray*}
\sig_i\sig_{j}\sig_i &=& \sig_{j}\sig_i\sig_{j} \quad |i-j| = 1 \\
\sig_i \sig_j &=& \sig_j\sig_i \quad\quad |i-j| \geq 2
\end{eqnarray*}

for instance, the braids
$\sig_1,\sig_2^{-1}\in B_3$ are given by
\[
\sig_1 = \braid{.5}{\r\s} \quad\quad\text{ and }\quad\quad
\sig_2^{-1} = \braid{.5}{\s\l}
\]

the group product then corresponds to vertical concatenation of
braids, so that
\[
\sig_2^{-1}\sig_1 = \braid{1}{\r\s\N\s\l}
\]

\end{slide}

%%%%%%%%%%%%%%%%%%%%%%%%%%%%%%%%%%%%%%%%%%%%%%%%%%%%%%%%%%%%%%%%%%%%%%%%%

\begin{slide}{Links as plat closures of braids}

\[
%b = \sigma_2 \sigma_3^{-1}\sigma_2 \;\; = \;\;
\braid{1.5}{
\s\r\s\N
\s\s\l\N
\s\r\s
} \stackrel{\text{plat}}{\longrightarrow}
\braid{1.5}{
\caps{0}{1}\caps{2}{1}\s\r\s\N
\s\s\l\N
\s\r\s\N
\cups{0}{1}\cups{2}{1}
} \,\,\,\rightsquigarrow\,
% \]
% The plat closure of $b$ is isotopic to the trefoil knot, as 
% \[
\xygraph{
!{0;/r1.0pc/:}
!P3"a"{~>{}}
!P9"b"{~:{(1.3288,0):}~>{}}
!P3"c":{~:{(2.5,0):}~>{}}
!{\vunder~{"b2"}{"b1"}{"a1"}{"a3"}}
!{\vcap~{"c1"}{"c1"}{"b4"}{"b2"}}
!{\vunder~{"b5"}{"b4"}{"a2"}{"a1"}}
!{\vcap~{"c2"}{"c2"}{"b7"}{"b5"}}
!{\vunder~{"b8"}{"b7"}{"a3"}{"a2"}}
!{\vcap~{"c3"}{"c3"}{"b1"}{"b8"}}
}
\]

%%%%%%%%%%%%%%%%%%%%%%%%%%%%%%%%%%%%%%%%%%%%%%%%%%%%%%%%%%%%%%%

\end{slide}

\begin{slide}{Birman moves}

two braids $b$ and $b'$ have \textcolor{blue}{isotopic closures} iff we
can go from $b$ to $b'$ by applying a sequence of
\textcolor{blue}{Birman moves}

% operations for the first Birman move
\[
\sig_{2i}\sig_{2i+1}\sig_{2i-1}\sig_{2i} =\! 
\braid{1.5}{\s\r\s\N\r\r\N\s\r\s},\,
\sig_{2i}\sig_{2i+1}\sig_{2i-1}^{-1}\sig_{2i}^{-1} =\! 
\braid{1.5}{\s\l\s\N\r\l\N\s\r\s},\,
\sig_{2i-1}^{\pm 1}
\]

\bigskip
\textcolor{green}{first Birman move:} $b\mapsto u b v$, where
$u,v$ are elements of the Birman subgroup generated by the above braids

\bigskip
\textcolor{green}{second Birman move:} $b\mapsto \sig_{2m}^{\pm 1}b \in
B_{2m+2}$, where two strands are added to the right of $b$ before
multiplying with $\sig_{2m}$

\end{slide}

\begin{slide}{References}
\begin{itemize}
\item V.~F.~R.~Jones, ``Hecke algebra representations of braid groups and link polynomials,'' \emph{Annals of Mathematics}, vol.~126, pp.~335-388, 1987.
\item H.~Wenzl, ``Hecke algebras of type $A_n$ and factors,'' \emph{Inventiones Mathematicae}, vol.~92, pp.~349-383, 1988.
\end{itemize}
\end{slide}

\begin{slide}{Jones polynomial}
%
% L_+ L_- L_0
%
\[
\xy
(6,6)*{}="tl";
(-6,6)*{}="tr";
(6,-6)*{}="bl";
(-6,-6)*{}="br";
{\ar|{\hole \; } "bl";"tr"};
{\ar "br";"tl"};
(0,-10)*{L_{+}};
\endxy
\qquad \qquad
\xy
(-6,6)*{}="tl";
(6,6)*{}="tr";
(-6,-6)*{}="bl";
(6,-6)*{}="br";
{\ar|{\hole \; } "bl";"tr"};
{\ar "br";"tl"};
(0,-10)*{L_{-}};
\endxy
\qquad\qquad 
 \xy
(-6,6)*{}="tl";
(6,6)*{}="tr";
(-6,-6)*{}="bl";
(6,-6)*{}="br";
"bl";"tl" **\crv{(-1,0)}?(1)*\dir{>};
"br";"tr" **\crv{(1,0)}?(1)*\dir{>};
(0,-10)*{L_{0}};
\endxy
\]
denote three oriented links that differ in a small region, as
symbolized by the above diagrams, but are otherwise the same

\medskip
the \textcolor{green}{Jones polynomial} $J_L(t)$ is a invariant of
oriented links that satisfies the skein relation
\[
t^{-1} J_{L_{+}}(t) - t J_{L_{-}}(t) = 
(t^{1/2} - t^{-1/2}) J_{L_{0}}(t)
\]

%%%%%%%%%%%%%%%%%%%%%%%%%%%%%%%%%%%%%%%%%%%%%%%%%%%%%%%%%%%%%%%%%%%%%%%%%

\end{slide}

\begin{slide}{Computation of the Jones polynomial}

with the help of the Kauffman bracket
\[
\left\langle \hspace{-.18in} \rule{0in}{.2in}
\xy
(3,3)*{}="tl";
(-3,3)*{}="tr";
(3,-3)*{}="bl";
(-3,-3)*{}="br";
{\ar@{-}|{\hole \; \hole \; \hole \; \hole \; \hole \; \hole } "bl";"tr"};
{\ar@{-} "br";"tl"};
%(0,-10)*{L_{\rm vert}};
\endxy \hspace{-.18in}
\right\rangle = A
\left\langle \, \rule{0in}{.2in}
\xy
(-3,3)*{}="tl";
(3,3)*{}="tr";
(-3,-3)*{}="bl";
(3,-3)*{}="br";
"bl";"tl" **\crv{(-1,0)};
"br";"tr" **\crv{(1,0)};
%(0,-10)*{L_{\rm cross}};
\endxy \,
\right\rangle + A^{-1}
\left\langle\, \rule{0in}{.2in}
\xy
(-3,3)*{}="tl";
(3,3)*{}="tr";
(-3,-3)*{}="bl";
(3,-3)*{}="br";
"tl";"tr" **\crv{(0,1)};
"bl";"br" **\crv{(0,-1)};
%(0,-10)*{L_{{\rm hor}}};
\endxy\,
\right\rangle
\]

the Jones polynomial is
\[
J_{L}(t) =  
A^{3w(\overrightarrow{L})}\left\langle L\right\rangle\Big|_{A = t^{1/4}}
\]
where $w(L)$ is the writhe of $L$

\medskip
connection between quantum computing and the Jones polynomial is based on

\medskip
\textcolor{red}{the Jones polynomial of plat closures evaluated at
roots of unity can be expressed with the help of a certain irrep of
$B_n$}

\end{slide}

\begin{slide}{Hecke algebra $H_n(q)$ - deformation of $\mathbb{C} S_n$} 

Hecke algebra $H_n(q)$ is generated by $g_1,\dotsc,g_{n-1}$
satisfying the relations
\begin{eqnarray*}
g_i^2 &=& g_i(q-1) + q \\
g_i g_{j} g_i &=&  g_{j}g_ig_{j} \quad\quad\quad\quad\; |i-j|=1 \\
g_ig_j  &=&  g_j g_i \quad\quad\quad\quad\quad\,  |i-j| > 1 
\end{eqnarray*}
this is a deformation of the symmetric group, which is
obtained when $q =1$

\medskip
we are interested in the case $q=e^{2\pi i/l}$, $q$ is $l$th primitive
root of unity 

\medskip
the map $B_n \rightarrow H_n(q)$ defined by $\sigma_i\mapsto g_i$ is a
representation of the braid group

\end{slide}

\begin{slide}{Quantum integers}

define the quantum integers
\[
{[d]}_\ell = \left.\frac{q^{d/2} - q^{-d/2}}{q^{1/2} - q^{-1/2}}
\right|_{q=e^{2\pi i/\ell}} = \frac{\sin(\pi d/\ell)}{\sin(\pi/\ell)}
\]

\bigskip
the loop value is
\[
D \equiv {[2]}_\ell = 
\left.q^{1/2} + q^{-1/2}\right|_{q=e^{2\pi i/\ell}} = 2\cos(\pi/\ell)
\]
\end{slide}

\begin{slide}{$H_n(q)$ in terms of projectors}

define projectors $e_i = (q-g_i)/(1+q)$ by setting
\[
g_i = q(1-e_i) - e_i = q - (1+q) e_i
\]

\bigskip
using these projectors, the relations of $H_n(q)$ may be expressed as
\begin{eqnarray*}
e_i^2 & = & e_i  \\
e_i e_j e_i - \tau e_j & = & e_{j}e_ie_{j} - \tau e_i \quad\quad |i-j| = 1\\
e_i e_j & = & e_j e_i \quad\quad\quad\quad\quad\,\, |i-j| > 1
\end{eqnarray*}
where $\tau=D^{-2}$
\end{slide}

%%%%%%%%%%%%%%%%%%%%%%%%%%%%%%%%%%%%%%%%%%%%%%%%%%%%%%%%%%%%%%%%%%%%%%%

\begin{slide}{Young lattice $\Lambda_n$}

\begin{tiny}
\[\xymatrix@R=.5cm@C=.3cm{
& & & & \emptyset \ar[d]\\
& & & & {\yng(1)} \ar[dr] \ar[dl]\\
& & & {\yng(2)}\ar[dl]\ar[dr] & &  {\yng(1,1)}\ar[dl]\ar[dr] \\
& & {\yng(3)}\ar[dll]\ar[d] & & {\yng(2,1)}\ar[dll]\ar[d]\ar[drr] & & {\yng(1,1,1)}\ar[d]\ar[drr] \\
{\yng(4)} & & {\yng(3,1)} & & {\yng(2,2)} & &  {\yng(2,1,1)} & & {\yng(1,1,1,1)}
}\]
\end{tiny}
\label{younggraph}

\end{slide}

%%%%%%%%%%%%%%%%%%%%%%%%%%%%%%%%%%%%%%%%%%%%%%%%%%%%%%%%%%%%%%%%%%%

\begin{slide}{Standard tableaux $T_n$} 

a \textcolor{green}{standard tableau} $T_\lambda$ is an assignment of
the numbers $1,2,\ldots,n$ to the boxes of $\lambda\in\Lambda_n$ such
that the numbers increase in each row and each column of $\lambda$

\bigskip
standard tableaux correspond to paths on the Young lattice
\end{slide}

%%%%%%%%%%%%%%%%%%%%%%%%%%%%%%%%%%%%%%%%%%%%%%%%%%%%%%%%%%%%%%%%%%%

\begin{slide}{truncated Young lattice $\Lambda_n^{(k,l)}$}
of $(k,l)$-admissible Young diagrams

\vspace{1.5cm}for example, $\Lambda_4^{(2,5)}$\quad\vspace{-3cm}
{\tiny
\[
\xymatrix@R=.5cm@C=.3cm{
& & & & \emptyset \ar[d]\\
& & & & {\yng(1)} \ar[dr] \ar[dl]\\
& & & {\yng(2)} \ar[dr] \ar[dl]& &  {\yng(1,1)}\ar[dl] \\
& & {\yng(3)} \ar[d]& & {\yng(2,1)}\ar[d]\ar[dll]  \\
& & {\yng(3,1)}  & & {\yng(2,2)}
}
\]
}

\[
\Lambda_n^{(k,\ell)} = \{\lambda\in \Lambda_n : 
\lambda_{k+1} = 0, \lambda_1 - \lambda_k\leq \ell-k\}
\]
\end{slide}

%%%%%%%%%%%%%%%%%%%%%%%%%%%%%%%%%%%%%%%%%%%%%%%%%%%%%%%%%%%%%%%%%%%%%%%%%

\begin{slide}{$(k,\ell)$-admissible standard tableaux $T_n^{(k,\ell)}$}

correspond to paths on $\Lambda_n^{(k,\ell)}$
\end{slide}

%%%%%%%%%%%%%%%%%%%%%%%%%%%%%%%%%%%%%%%%%%%%%%%%%%%%%%%%%%%%%%%%%%%%%%%%

\begin{slide}{Jones-Wenzl irreps of $H_n(q)$ and $B_n$}
enumerated by $(k,\ell)$-admissible Young diagrams
$\lambda\in\Lambda_n^{(k,\ell)}$

\bigskip
corresponding \textcolor{green}{irrep $\pi_\lambda^{(k,l)}$} acts on
the vector space
\[
\mbox{\textcolor{green}{$V_\lambda^{(k,\ell)}$}} = 
{\rm span}\{ \ket{t}\, :\, t\in T_n^{(k,\ell)}\}
\]
whose basis vectors are the $(k,\ell)$-admissible standard tableaux of
shape $\lambda$

\bigskip
define \textcolor{green}{$s_i(t)$} to be the standard tableau obtained
form $t$ by exchanging the positions of $i$ and $i+1$ (not valid, set
$s_i(t)=t$)

\bigskip
define \textcolor{green}{hook length} $d_i(t)=c_i(t)-c_{i+1}(t) -
(r_i(t)-r_{i+1}(t))$

\end{slide}

%%%%%%%%%%%%%%%%%%%%%%%%%%%%%%%%%%%%%%%%%%%%%%%%%%%%%%%%%%%%%%%%%%%%%%%%


\begin{slide}{Jones-Wenzl irreps continued}

let $V_{t,i}^{(k,l)}$ be the subspace of $V_\lambda^{(k,\ell)}$ spanned by
$\ket{t}$ and $\ket{s_i(t)}$
\[
\hspace{-7.1cm}\mbox{define}\quad
a_\ell(d) = \frac{[d+1]_\ell}{[2]_\ell [d]_\ell}
\]

\medskip
the \textcolor{green}{image $\pi_\lambda^{(k,\ell)}(e_i)$} of $e_i$ under
$\pi_\lambda^{(k,\ell)}$ acts as
\[
\pi_\lambda^{(k,\ell)}(e_i) = 
\pmat{a_\ell(d_i(t)) & \sqrt{a_\ell(d_i(t))(1-a_\ell(d_i(t)))} \\
\sqrt{a_\ell(d_i(t))(1-a_\ell(d_i(t)))} & 1-a_\ell(d_(i(t))}
\]

\medskip
on $V_{t,i}^{(k,\ell)}$ and as identity on the orthogonal complement
\end{slide}

%%%%%%%%%%%%%%%%%%%%%%%%%%%%%%%%%%%%%%%%%%%%%%%%%%%%%%%%%%%%%%%%%%%%%%%%%

\begin{slide}{Jones-Wenzl irreps continued}

corresponding \textcolor{green}{irrep of the braid group $B_n$} is
defined by
\[
\pi_\lambda^{(k,\ell)}(g_i) =
q\, {\bf 1}_{V_\lambda^{(k,l)}} -(1+q) \pi_\lambda^{(k,\ell)}(e_i)
\]
\end{slide}

%%%%%%%%%%%%%%%%%%%%%%%%%%%%%%%%%%%%%%%%%%%%%%%%%%%%%%%%%%%%%%%%%%%%%%%%%

\begin{slide}{Representation-theoretic formula} 
let $\lambda=[m,m]$ be the ``rectangular'' Young diagram with two
rows, each having $m$ boxes

\newcommand{\nnn}{\mbox{\tiny $2m\!\!-\!\!1$}}
\newcommand{\nnnn}{\mbox{\tiny $2m$}}

\medskip
define the normalized vector
\[
\ket{\varphi_{2m}} = \ket{\,
{\Yvcentermath1\Yboxdim22pt\young(135,246)}\cdots
{\Yvcentermath1\Yboxdim22pt\young(\nnn,\nnnn)}\,}
\]
\end{slide}

\begin{slide}{Representation-theoretic formula}
\textcolor{green}{the norm of the Jones polynomial of the plat
closure of any $b\in B_{2m}$ evaluated at $e^{2\pi i /l}$ is given by}
\[D^{m-1}
\left|
\bra{\varphi_{2m}} \pi_\lambda^{(2,l)}(b) \ket{\varphi_{2m}}
\right|
\]

\medskip
\textcolor{red}{has the same structure as $|\bra{00\ldots 0} U
\ket{00\ldots 0}|$ in the BQP-complete problem ``Approximate circuit
matrix entry''} 
\end{slide}

%%%%%%%%%%%%%%%%%%%%%%%%%%%%%%%%%%%%%%%%%%%%%%%%%%%%%%%%%%%%%%%%%%%%%%%%%%%%

%%%%%%%%%%%%%%%%%%%%%%%%%%%%%%%%%%%%%%%%%%%%%%%%%%%%%%%%%%%%%%%%%%%%%%%%%

\begin{slide}{Our encoding of a qubit in $4$ strands}

\[
\xymatrix@R=.1cm{
& & {\young(1,2)} \ar[rd] \\
 & {\yng(1)} \ar[ur] \ar[dr] & & {\yng(2,1)} \ar[r] & 
{\yng(2,2)} \\
& & {\young(12)} \ar[ru]
}
\]

\medskip
\[
\ket{0} \equiv \ket{\, \Yvcentermath1 \young(13,24)\,}  \quad\quad
\ket{1} \equiv \ket{\, \Yvcentermath1 \young(12,34)\,}
\]

\end{slide}

%%%%%%%%%%%%%%%%%%%%%%%%%%%%%%%%%%%%%%%%%%%%%%%%%%%%%%%%%%%%%%%%%%%%%%%%%

\begin{slide}{Our encoding of $2$ qubits in $8$ strands}

\[
\!\!\!\!
\xymatrix@R=.01cm{
                 &                           & {\yng(1,1)} \ar[rd] \\
 & {\yng(1)} \ar[ur] \ar[dr] &                     &
	  {\yng(2,1)} \ar[r] & {\yng(2,2)} \ar[r] &  \\
& & {\yng(2)} \ar[ru]
}
\]

\[
\xymatrix@R=.01cm{
                 &                           & {\yng(3,3)} \ar[rd] \\
 & {\yng(3,2)} \ar[ur] \ar[dr] &                     & {\yng(4,3)} \ar[r] & {\yng(4,4)} \\
& & {\yng(4,2)} \ar[ru]
}
\]

\end{slide}

%%%%%%%%%%%%%%%%%%%%%%%%%%%%%%%%%%%%%%%%%%%%%%%%%%%%%%%%%%%%%%%%%%%%%%%%%

\begin{slide}{Encoding of $2$ qubits into $8$ strands}

four periodic structure
\vspace{-0.5cm}
\begin{eqnarray*}
\ket{t_0}\ket{t_0} &\mapsto&
  \ket{\,\Yvcentermath1 \young(1357,2468)\,} \\
\ket{t_0}\ket{t_1} &\mapsto&
  \ket{\,\Yvcentermath1 \young(1356,2478)\,} \\
\ket{t_1}\ket{t_0} &\mapsto&
  \ket{\,\Yvcentermath1 \young(1257,3468)\,} \\
\ket{t_1}\ket{t_1} &\mapsto&
  \ket{\,\Yvcentermath1 \young(1256,3478)\,}
\end{eqnarray*}

\medskip
${\cal
C}^{(2)}=\text{span}\{\ket{t_0}\ket{t_0},\ldots,\ket{t_1}\ket{t_1}\}$
is a subspace of $V_{\Yboxdim{3pt}\yng(4,4)}$ 

\medskip
${\cal C}^{(n)}=\text{span}\{\ket{t_0},\ket{t_1}\}^{\otimes n} \le 
V_{\Yvcentermath1\Yboxdim{3pt}\yng(2,2) \cdots \Yvcentermath1\Yboxdim{3pt}\yng(1,1)}$

\end{slide}

%%%%%%%%%%%%%%%%%%%%%%%%%%%%%%%%%%%%%%%%%%%%%%%%%%%%%%%%%%%%%%%%%%%%%%%%%

\begin{slide}{Our encoding of two-qubit gates}

makes use of 

\bigskip
\textbf{Theorem} [Freedman, Larsen, Wang] 

\bigskip
the image of the eight strand braid group $B_{8}$ under the
irreducible representation $\pi_{\Yboxdim{3pt}\yng(4,4)}{}$ is dense
in the special unitary group $SU\big(V_{\Yboxdim{3pt}\yng(4,4)}\big)$

\begin{center}{\large $\Downarrow$}\end{center}

\medskip
\textcolor{green}{any two-qubit gate $U$ can be efficiently
approximated by sequences elementary operations $\sigma_1^{\pm
1},\ldots,\sigma_7^{\pm 1}$}

\bigskip
\textcolor{blue}{such sequences can be efficiently found by a
classical computer} (Solovay-Kitaev theorem)

\end{slide}

%%%%%%%%%%%%%%%%%%%%%%%%%%%%%%%%%%%%%%%%%%%%%%%%%%%%%%%%%%%%%%%%%%%%%%%%%

\begin{slide}{Efficient encoding of quantum circuits}

\textbf{Theorem} [W. and Yard]

\medskip
let $U=U_{m} U_{m-1} \cdots U_1$ be a quantum circuit of length $m$
acting on $n$ qubits and let $\epsilon > 0$

\medskip
then there is a braid $b\in B_{4n}$ with $poly(m, log(1/\epsilon))$
crossings such that

\[
\|\, U\, -\, \pi(b)|_{{\cal C}^{(n)}}\, \|_{\infty}
\leq
\epsilon
\]

\medskip
$\pi(b)|_{{\cal C}^{(n)}}$ denotes restriction to computational subspace

\bigskip
moreover, such a braid can be found from a description of the circuit in
$poly(m, log(1/\epsilon))$ time on a classical computer

\end{slide}

%%%%%%%%%%%%%%%%%%%%%%%%%%%%%%%%%%%%%%%%%%%%%%%%%%%%%%%%%%%%%%%%%%%%%%%%%

\begin{slide}{Putting everything together}
\begin{itemize}
\item let ${\cal G}$ be any universal set and $\epsilon=1/poly(n)$
\item let $U$ be a quantum circuit acting on $n$ qubits and using $m=poly(n)$ gates from ${\cal G}$ 
\end{itemize}
given $U$ and $\epsilon$ we can efficiently construct a braid $b\in B_{4n}$ of
length $poly(n)$ such that

\[
\left|\, 
|\bra{00\ldots 0} U \ket{00\ldots 0}| - 
|\bra{\varphi} \pi^{(2,l)}(b) \ket{\varphi}| \,
\right| < \epsilon/2
\]
\[
\left|\, 
|\bra{00\ldots 0} U \ket{00\ldots 0}| - 
D^{-2n+1}| J_{\hat{b}}(e^{2\pi i/l})|\, \right| < \epsilon/2
\]
where $J_{\hat{b}}(e^{2\pi i/l})$ is the Jones polynomial of the plat
closure of $b$ evaluated at the root of unity $t=e^{2\pi i/l}$
\end{slide}

%%%%%%%%%%%%%%%%%%%%%%%%%%%%%%%%%%%%%%%%%%%%%%%%%%%%%%%%%%%%%%%%%%%%%%

\begin{slide}{Putting every thing together continued}

if we can estimate $|J_{\hat{b}}(e^{2\pi i/l})|$ with
precision $\epsilon \cdot D^{2n-1}/2$, then 

we can solve the BQP-complete problem ``approximate circuit
matrix entry''

\bigskip
this shows that approximately evaluating the Jones polynomial at roots of unity is BQP-hard
\end{slide}

%%%%%%%%%%%%%%%%%%%%%%%%%%%%%%%%%%%%%%%%%%%%%%%%%%%%%%%%%%%%%%%%%%%%%%

\begin{slide}{More results}
\begin{itemize}
\item Approximately evaluating the Jones polynomial of the plat and trace closure of braids are in BQP
\item Approximately evaluating the HOMFLYPY polynomial of trace closure of braids at certain pairs of values is in BQP
\item Simplified and new proofs of the computational complexity of certain problems concerning the Jones polynomial 
\end{itemize}
\end{slide}

\begin{slide}{Historical overview}
\begin{itemize}
\item ``Approximately evaluating the Jones polynomial is BQP-complete'' proved in the context of topological quantum computation; uses topology and quantum field theories [Freedman, Kitaev, Larsen, Wang]
\item ``Approximately evaluating the Jones polynomial is in BQP''; simplified proof within the quantum circuit model [Aharonov, Jones, Zhang]
\item ``Approximately evaluating the Jones polynomial is in BQP-hard''; simplified proof based on a certain encoding of quantum circuits in the rectangular representation of the braid group [W. and Yard]
\end{itemize}
\end{slide}

\end{document}
















































\begin{slide}{Single qubits}
\textcolor{green}{classical bit (bit):}\, $\mathbb{B}:=\{0,1\}$
\begin{itemize}
\item state: either $0$ or $1$
\end{itemize}

\bigskip
\textcolor{green}{quantum bit (qubit):}\, $\mathbb{C}^2$
\begin{itemize}
\item basis states (classical states):
\[
|0\>:=\left(\begin{array}{c}1\\0\end{array}\right)
\quad\quad\mbox{and}\quad\quad
|1\>:=\left(\begin{array}{c}0\\1\end{array}\right)
\]
\item general states: \textcolor{red}{superpositions}
\[
|\psi\> = 
\left(
\begin{array}{c}
\alpha\\
\beta
\end{array}
\right)=
\alpha |0\> + \beta |1\>\,,\quad |\alpha|^2+|\beta|^2=1
\]
\end{itemize}
\end{slide}

\begin{slide}{Quantum register}
\textcolor{green}{Quantum register:}
\[
\underbrace{\mathbb{C}^2\otimes\mathbb{C}^2\otimes\cdots\otimes\mathbb{C}^2}_{n}
\]

$2^n$-dimensional vector space with \textcolor{red}{tensor product structure}

\bigskip
\textcolor{green}{general state:} 
\[
|\psi\>=\sum_{x\in\{0,1\}^n} \alpha_x |x\>
\]
where $\sum_{x} |\alpha_x|^2 = 1$ and
$|x\>=|x_1\>\otimes|x_2\>\otimes\cdots\otimes|x_n\>$

\medskip

\bigskip
\textcolor{red}{exponentially many basis states in superposition!}
\end{slide}

\begin{slide}{Single qubit gates on two qubits}
one-qubit operation $U$ on \textcolor{blue}{second} qubit

{
\Qcircuit @C=1em @R=1em {
& \qw      & \qw \\
& \gate{{\blue U}} & \qw 
}}

\medskip
action on basis states of $\mathbb{C}^2\otimes\mathbb{C}^2$\quad
$
\begin{array}{ccc}
|00\> & \mapsto & |0\>\, {\blue U}|0\>\\
|01\> & \mapsto & |0\>\, {\blue U}|1\>\\
|10\> & \mapsto & |1\>\, {\blue U}|0\>\\
|11\> & \mapsto & |1\>\, {\blue U}|1\>
\end{array}
$

corresponding matrix

$ {\red I}\otimes {\blue U}
=
\left(
\begin{array}{c|c}
{\red 1} \cdot {\blue U} & {\red 0} \cdot {\blue U} \\ \hline
{\red 0} \cdot {\blue U} & {\red 1} \cdot {\blue U}
\end{array}
\right)
=
\left(
\begin{array}{cccc}
{\blue u_{00}} & {\blue u_{01}} & 0      & 0 \\
{\blue u_{10}} & {\blue u_{11}} & 0      & 0 \\
0              & 0              & {\blue u_{00}} & {\blue u_{01}} \\
0              & 0              & {\blue u_{10}} & {\blue u_{11}}
\end{array}
\right)
$
\end{slide}

\begin{slide}{Single qubit gates on two qubits}
{
\Qcircuit @C=1em @R=1em {
& \gate{{\red U}} & \qw \\
& \qw      & \qw 
}}

action on basis states of $\mathbb{C}^2\otimes\mathbb{C}^2$\quad
$
\begin{array}{ccc}
|00\> & \mapsto & ({\red U}|0\>)\, |0\>\\
|01\> & \mapsto & ({\red U}|0\>)\, |1\>\\
|10\> & \mapsto & ({\red U}|1\>)\, |0\>\\
|11\> & \mapsto & ({\red U}|1\>)\, |1\>
\end{array}
$

\medskip
corresponding matrix 

${\red U}\otimes {\blue I}
=
\left(
\begin{array}{c|c}
{\red u_{00}} \cdot {\blue I} & {\red u_{01}} \cdot {\blue I} \\ \hline
{\red u_{10}} \cdot {\blue I} & {\red u_{11}} \cdot {\blue I}
\end{array}
\right)
=
\left(
\begin{array}{cccc}
{\red u_{00}} & 0             & {\red u_{01}} & 0      \\
0             & {\red u_{00}} & 0             & {\red u_{01}} \\
{\red u_{10}} & 0             & {\red u_{11}} & 0      \\
0             & {\red u_{10}} & 0             & {\red u_{11}}
\end{array}
\right)
$
\end{slide}


%%%%%%%%%%%%%%%%%%%%%%%%%%%%%%%%%%%%%%%%%%%%%%%%%%%%%%%%%%%%%%%%%%%%%%%%%%%%s

\begin{slide}{Controlled CNOT gate}

{
\Qcircuit @C=1em @R=1em {
& \ctrl{1} & \qw \\
& \targ     & \qw 
}}

action on basis states of $\mathbb{C}^2\otimes\mathbb{C}^2$

\[
\textrm{CNOT}|a\>|b\> = |a\>|a\oplus b\>
\]

corresponding matrix
\[
\left(
\begin{array}{cccc}
1      & 0      & 0      & 0 \\
0      & 1      & 0      & 0 \\
0      & 0      & 0      & 1 \\
0      & 0      & 1      & 0
\end{array}
\right)
\]

\end{slide}

\begin{slide}{Decomposition into elementary gates}

Theorem: 

\medskip
unitary operations $U\in {\cal U}(2^n)$ on $n$ qubits can be
implemented by quantum circuits using only

\begin{itemize}
\item \textcolor{blue}{CNOT} gates 

\medskip
\item \textcolor{blue}{single qubit} gates
\end{itemize}
\end{slide}

%%%%%%%%%%%%%%%%%%%%%%%%%%%%%%%%%%%%%%%%%%%%%%%%%%%%%%%%%%%%%%%%%%%%%%%%%%%%%%%%%%%%%%%%%%%%%%%%%%%

\begin{slide}{Decomposition of the Tofolli gate}
\Qcircuit @C=0.7em @R=0.7em {
& \qw      & \qw      & \qw      & \ctrl{1}      & \ctrl{2} & \qw           & \ctrl{1} & \qw      & \gate{W}      & \ctrl{2} & \qw 
& \qw \\
%
& \qw      & \gate{W} & \ctrl{1} & \targ         & \qw      & \gate{W^\dag} & \targ    & \ctrl{1} & \qw           & \qw      & \qw 
& \qw \\
%
& \gate{H} & \gate{W} & \targ    & \gate{W^\dag} & \targ    & \gate{W}      & \qw      & \targ    & \gate{W^\dag} & \targ    & \gate{H}
& \qw
}

\medskip
$W={\rm diag}(1,e^{i\pi/4})$ \hspace{1cm} $H=
\frac{1}{\sqrt{2}}
\left(
\begin{array}{cc}
1 & 1 \\
1 & -1
\end{array}
\right)
$

\end{slide}

\end{document}